\PassOptionsToPackage{usenames,dvipsnames}{xcolor}
\documentclass[a4paper,12pt]{letter}
\usepackage[utf8]{inputenc}
\usepackage[T1]{fontenc}
\usepackage[english,slovene]{babel}
\usepackage[tikz]{bclogo}
\usepackage{qrcode,marvosym}
\usepackage{pgf,tikz,pgfplots}
\pgfplotsset{compat=1.14}
\usepackage{mathrsfs}
\usetikzlibrary{arrows}
\usepackage{graphicx}
\newcommand{\ds}{\displaystyle}
\usepackage{fancybox}
\usepackage{amsfonts}
\usepackage{amsmath}
\usepackage{wallpaper}
\usepackage{hyperref}
\usepackage{array}
\newcolumntype{M}[1]{>{\centering\arraybackslash}m{#1}}
\newcolumntype{N}{@{}m{0pt}@{}}
\usepackage{fancyhdr,enumerate}
\newcommand{\ora}{\overrightarrow}
\renewcommand{\headrulewidth}{3pt}
\renewcommand{\headrule}{\hbox to\headwidth{%
  \color{gray}\leaders\hrule height \headrulewidth\hfill}}
\renewcommand{\footrulewidth}{2pt}
\renewcommand{\footrule}{\hbox to\headwidth{%
  \color{brown}\leaders\hrule height \footrulewidth\hfill}}
\usepackage{gensymb}
\usepackage{amssymb}
  \textwidth 20 true cm
  \oddsidemargin -2 true cm
  \evensidemargin -2 true cm
  \topmargin -3.5 true cm
  \textheight 27.5 true cm
  \linespread{1.2}
\renewcommand{\headrulewidth}{0.3pt}
\renewcommand{\footrulewidth}{1.9pt}
\definecolor{qqwuqq}{rgb}{1,0,0}
\definecolor{qqqqff}{rgb}{0,0,1}
\definecolor{qqffqq}{rgb}{0,1,0}
\definecolor{ffqqqq}{rgb}{1,0,0}
\definecolor{zzuxux}{rgb}{0,0,0}
\usepackage{mdframed}
\usepackage{multicol}
\definecolor{bgblue}{RGB}{0,0,253}
\usepackage{eurosym}
%%%%%%%%%%%%%%%%%%%%%%%%%%%%%%%%%%%%%%%%%%%%%%%%%%%%%%%%
\newcounter{tocke}
\setcounter{tocke}{0}
\usepackage{tikz-3dplot}
%%%%%%%%%%%%%%%%%%%%%%%%%%%%%%%%%%%%%%%%%%%%%%%%%%%%%%%%%%
\mdfdefinestyle{theoremstyle}{%
linecolor=brown!40,linewidth=1pt,%
frametitlerule=true,%
frametitlebackgroundcolor=brown!50,
innertopmargin=\topskip,
}
\mdfdefinestyle{theoremstyle1}{%
linecolor=brown,middlelinewidth=1pt,%
frametitlerule=true,%
apptotikzsetting={\tikzset{mdfframetitlebackground/.append style={%
shade,left color=black!40, right color=gray!10},
mdfbackground/.append style={
top color=white!100!white,
bottom color=black!5
},
}},
frametitlerulecolor=gray,
frametitlerulewidth=2pt,
innertopmargin=\topskip,
innerbottommargin=\topskip,
}
\mdtheorem[style=theoremstyle1]{naloga}{Naloga}
\mdtheorem[style=theoremstyle]{resitev}{Analiza Naloge}
\mdtheorem[style=theoremstyle]{kriterij}{\v{S}tevilo dose\v{z}enih to\v{c}k na testu}
%%%%%%%%%%%%%%%%%%%%%%%%%%%%%%%%%%%%%%%%%%%%%%%%%%%%%%%%%%
\usepackage[table]{xcolor}
\usepackage[printwatermark]{xwatermark}
\newwatermark[allpages,color=brown!2,angle=45,scale=5,
xpos=-5,ypos=0]{matej.info}
%%%%%%%%%%%%%%%%%%%%%%%%%%%%%%%%%%%%%%%%%%%%%%%%%%%%%%%%%%

\begin{document}
\pagestyle{fancy}
\renewcommand\headrulewidth{0pt}
\lhead{}\chead{}\rhead{}
\rfoot{\vspace*{-2\baselineskip}}
\renewcommand{\vec}{\overrightarrow}
%%%%%%%%%%%%%%%%%%%% podatki testa -- tu spreminjaj
\def\programletnik{Gim - 3}
\def\naslov{Test 1.1 - }
\def\sklop{Eksponentna in logaritemska funkcija}
%%%%%%%%%%%%%%%%%%%% tocke po nalogah
\def\tockeA{5+5}
\def\tockeB{5}
\def\tockeC{4+4}
\def\tockeD{4}
\def\tockeE{4+4}
\def\tockeF{5}
%%%%%%%%%%%%%%%%%%%%
\begin{minipage}{20cm}
\begin{bclogo}[couleur=brown!30,
  arrondi=0,ombre=true, couleurOmbre=gray!50,
logo =\bcplume,
  noborder=true]
{\textrm{{\Large{\naslov \sklop} \Huge{}
\large{} \hfill $\programletnik$}}}
\end{bclogo}
\end{minipage}
\begin{minipage}{20cm}
\begin{bclogo}[couleur=brown!30,
  arrondi=0,ombre=true, couleurOmbre=gray!50,
logo =\bcplume,
  noborder=true]
{\textrm{{\raisebox{-0.3cm}{\large{Ime in Priimek: \rule{10cm}{0.1mm}} \Huge{}
\large{} \hfill $ $}}}}
\end{bclogo}
\end{minipage}
%%%%%%%%%%%%%%%%%%%%%%%%%%%%%%%%%%%%%%%%%%%%%%%%%

%%%%%%%%%%%%%%%%%%%%%%%%%%%%%%%%%%%%%%%%%%%%%%%%%
\begin{naloga}[\hfill $\tockeA$] \addtocounter{tocke}{\numexpr\tockeA\relax}
Dana je funkcija $f(x)=\left(\dfrac12\right)^{x-2}-4$.
\begin{enumerate}[a)]
\item Izra\v{c}unaj ni\v{c}lo funkcije $f$ in dolo\v{c}i to\v{c}ko na grafu  $T\!\left(x,-\dfrac{127}{32}\right)$. Nari\v{s}i graf funkcije $f$ (ozna\v{c}i asimptoto ter obe dobljeni to\v{c}ki).
\item Zapi\v{s}i inverzno funkcijo $f^{-1}$. Dolo\v{c}i ena\v{c}bo asimptote in ni\v{c}lo funkcije $f^{-1}$.
\end{enumerate}

\end{naloga}

\vspace{0.3cm}
\begin{center}
\begin{tikzpicture}[scale=0.75]
\draw[very thin,gray!30] (-2,-5) grid (9,5);
\draw[->,thick] (-2,0) -- (9,0) node[right] {$x$};
\draw[->,thick] (0,-5) -- (0,5) node[above] {$y$};
\end{tikzpicture}
\newpage

%%%%%%%%%%%%%%%%%%%%%%%%%%%%%%%%%%%%%%%%%%%%%%%%%
\begin{naloga}[\hfill $\tockeB$] \addtocounter{tocke}{\numexpr\tockeB\relax}
Re\v{s}i ena\v{c}bo, rešitev zapiši na štiri mesta:
\[
3^{1-x}=2^{x}
\]
\end{naloga}

\vspace{8cm}

%%%%%%%%%%%%%%%%%%%%%%%%%%%%%%%%%%%%%%%%%%%%%%%%%
\begin{naloga}[\hfill $\tockeC$] \addtocounter{tocke}{\numexpr\tockeC\relax}
Re\v{s}i ena\v{c}bi:
\begin{enumerate}[a)]
\item $\log_3\big(\log_2(x+1)\big)=e^{\ln 2}$
\item $\log_2 x+\log_4 x+\dfrac32=0$
\end{enumerate}
\end{naloga}\newpage

%%%%%%%%%%%%%%%%%%%%%%%%%%%%%%%%%%%%%%%%%%%%%%%%%
\begin{naloga}[\hfill $\tockeD$] \addtocounter{tocke}{\numexpr\tockeD\relax}
Zapi\v{s}i definicijsko obmo\v{c}je funkcije
\[
f(x)=\ln\big((2x-1)(x+1)\big).
\]
\end{naloga}\vspace{8 cm}

%%%%%%%%%%%%%%%%%%%%%%%%%%%%%%%%%%%%%%%%%%%%%%%%%
\begin{naloga}[\hfill $\tockeE$] \addtocounter{tocke}{\numexpr\tockeE\relax}
\begin{enumerate}[a)]
\item Naj bo $\ln x=\dfrac13$, $\ln y=2$, $\ln z=-\dfrac14$. Izra\v{c}unaj
\[
\ln\left(\dfrac{e\,x^3}{\sqrt{y\,z^{-1}}}\right).
\]
\item Dolo\v{c}i $x$, \v{c}e je
\[
\log x=2\log 3+\dfrac12\log 9-\dfrac{1}{\log_{27}10}.
\]
\end{enumerate}
\end{naloga}\newpage

%%%%%%%%%%%%%%%%%%%%%%%%%%%%%%%%%%%%%%%%%%%%%%%%%
\begin{naloga}[\hfill $\tockeF$] \addtocounter{tocke}{\numexpr\tockeF\relax}
Izra\v{c}unaj produkt re\v{s}itev ena\v{c}be
\[
\log_2\!\left(x^{\log_2 x}\right)=\log_2\big(x^2\big)+\log_2 4.
\]
\end{naloga}

\vfill
\begin{kriterij*} \begin{center} \renewcommand{\arraystretch}{2.5}
\begin{tabular}{|>{\columncolor{gray!20}}c||c|c|c|c|c||c|c|c|} \hline Odstotki &$[0,45)$ & $[45,60)$ & $[60,75)$ & $[75,90)$ & $[90,100]$& \v{S}t. mo\v{z}nih to\v{c}k& Dose\v{z}ene to\v{c}ke & OCENA \\ \hline Ocena & 1 & 2 & 3 & 4 & 5 &\cellcolor{gray!20} $\thetocke$&\cellcolor{gray!20} \phantom{ocena} & \cellcolor{gray!20} \phantom{a} \\ \hline \end{tabular} \end{center} \vspace{0.3cm}  \end{kriterij*}

\newpage
\textsc{\Large Re\v{s}itve s postopki}
\vspace{0.5cm}

\begin{resitev}[-- eksponentna funkcija in njen inverz]
Dana je $f(x)=\left(\dfrac12\right)^{x-2}-4$. Funkcija je padajo\v{c}a, njena horizontalna asimptota je $y=-4$.

\textbf{a)} Ni\v{c}lo dobimo iz $f(x)=0$:
\[
\left(\frac12\right)^{x-2}=4=\left(\frac12\right)^{-2}\ \Rightarrow\ x-2=-2\ \Rightarrow\ x=0.
\]
Funkcija ima ni\v{c}lo pri $x=0$, torej gre graf skozi to\v{c}ko $(0,0)$.

Za drugo to\v{c}ko re\v{s}imo $f(x)=-\dfrac{127}{32}$:
\[
\left(\frac12\right)^{x-2}-4=-\frac{127}{32}\ \Rightarrow\ \left(\frac12\right)^{x-2}=4-\frac{127}{32}=\frac{128-127}{32}=\frac{1}{32}=\left(\frac12\right)^{5}.
\]
Sledi $x-2=5$, torej $x=7$ in $T\!\left(7,-\dfrac{127}{32}\right)$.

Graf je padajo\v{c}a eksponentna krivulja s horizontalno asimptoto $y=-4$, ki poteka skozi to\v{c}ki $(0,0)$ in $\left(7,-\dfrac{127}{32}\right)$ (slednja je prakti\v{c}no \v{z}e tik ob asimptoti).

\textbf{b)} Inverzno funkcijo dobimo tako, da zamenjamo vlogi $x$ in $y$ ter izrazimo novi $y$:
\[
y=\left(\frac12\right)^{x-2}-4\ \Longrightarrow\ x=\left(\frac12\right)^{y-2}-4\ \Longrightarrow\ x+4=\left(\frac12\right)^{y-2}
\]
\[
\Longrightarrow\ y-2=\log_{1/2}(x+4)\ \Longrightarrow\ y=\log_{1/2}(x+4)+2.
\]
Torej
\[
f^{-1}(x)=\log_{1/2}(x+4)+2,\qquad D_{f^{-1}}=(-4,\infty).
\]
Ker je graf $f^{-1}$ zrcalna slika grafa $f$ prek premice $y=x$, je asimptota funkcije $f^{-1}$ \textbf{navpi\v{c}na} premica
\[
x=-4.
\]
Ni\v{c}lo funkcije $f^{-1}$ dobimo iz $f^{-1}(x)=0$:
\[
\log_{1/2}(x+4)+2=0\ \Rightarrow\ \log_{1/2}(x+4)=-2\ \Rightarrow\ x+4=\left(\frac12\right)^{-2}=4\ \Rightarrow\ x=0.
\]
(To se ujema s tem, da je $f(0)=0$, zato mora biti tudi $f^{-1}(0)=0$.)

\begin{center}
\begin{tikzpicture}[scale=0.62]
\draw[very thin,gray!25] (-5,-5) grid (13,13);
\draw[->,thick] (-5,0) -- (13,0) node[right] {$x$};
\draw[->,thick] (0,-5) -- (0,13) node[above] {$y$};
\draw[dashed,thick,blue!60] (-5,-4) -- (13,-4) node[pos=0.03,below]{$y=-4$};
\draw[dashed,thick,red!60] (-4,-5) -- (-4,13) node[pos=0.97,right]{$x=-4$};
\draw[dotted] (-5,-5) -- (13,13);
\draw[domain=-1:8,smooth,variable=\x,blue,thick] plot ({\x},{0.5^(\x-2)-4});
\draw[domain=-3.9:12,smooth,variable=\x,red,thick,samples=100] plot ({\x},{-ln(\x+4)/ln(2)+2});
\filldraw[blue] (0,0) circle (2pt);
\filldraw[blue] (7,-3.969) circle (2pt) node[below right]{\scriptsize $T(7,-\frac{127}{32})$};
\filldraw[red] (-3.969,7) circle (2pt) node[above left]{\scriptsize $T'(-\frac{127}{32},7)$};
\node[blue] at (2.3,3.2) {$f$};
\node[red] at (3.2,2.3) {$f^{-1}$};
\end{tikzpicture}
\end{center}
\end{resitev}

\begin{resitev}[-- eksponentna ena\v{c}ba]
Ena\v{c}bo $3^{1-x}=2^x$ re\v{s}imo tako, da logaritmiramo obe strani (naravni logaritem):
\[
\ln\!\left(3^{1-x}\right)=\ln\!\left(2^{x}\right)\ \Longrightarrow\ (1-x)\ln3=x\ln2.
\]
Izpostavimo $x$:
\[
\ln3-x\ln3=x\ln2\ \Longrightarrow\ \ln3=x(\ln2+\ln3)=x\ln6.
\]
Sledi
\[
x=\frac{\ln3}{\ln6}=\log_6 3=1-\log_6 2\approx0{,}6131.
\]
\end{resitev}

\begin{resitev}[-- logaritemski ena\v{c}bi]
\textbf{a)} $e^{\ln 2} = 2$, $\log_3\big(\log_2(x+1)\big)=2$. Iz definicije logaritma sledi
\[
\log_2(x+1)=3^2=9\ \Longrightarrow\ x+1=2^9=512\ \Longrightarrow\ x=511.
\]
Preverimo definicijsko obmo\v{c}je: $x+1=512>0$ in $\log_2(x+1)=9>0$, zato je re\v{s}itev veljavna.

\textbf{b)} $\log_2 x+\log_4 x+\dfrac32=0$. Ker je $\log_4 x=\dfrac{\log_2 x}{\log_2 4}=\dfrac{\log_2 x}{2}$, uvedemo substitucijo $t=\log_2 x$:
\[
t+\frac t2+\frac32=0\ \Longrightarrow\ \frac32t=-\frac32\ \Longrightarrow\ t=-1.
\]
Torej $\log_2 x=-1$, od koder $x=\dfrac12$ (definicijsko obmo\v{c}je $x>0$ je izpolnjeno).
\end{resitev}

\newpage
\begin{resitev}[-- definicijsko obmo\v{c}je]
Ker je $\ln$ definiran le za pozitivne argumente, mora veljati
\[
(2x-1)(x+1)>0.
\]
Ni\v{c}li izraza sta $x=\dfrac12$ in $x=-1$. Ker je koeficient pri $x^2$ pozitiven ($2x^2+x-1$), je izraz pozitiven \textbf{zunaj} ni\v{c}el, torej za $x<-1$ ali $x>\dfrac12$. Definicijsko obmo\v{c}je je zato
\[
D_f=(-\infty,-1)\cup\left(\frac12,\infty\right).
\]
\end{resitev}

\begin{resitev}[-- ra\v{c}unanje z logaritmi]
\textbf{a)} Podano imamo $\ln x=\dfrac13$, $\ln y=2$, $\ln z=-\dfrac14$. Uporabimo pravila za logaritme (produkt, koli\v{c}nik, potenca, koren):
\[
\ln\left(\frac{e\,x^3}{\sqrt{y\,z^{-1}}}\right)=\ln e+3\ln x-\frac12\big(\ln y-\ln z\big).
\]
Vstavimo vrednosti:
\[
=1+3\cdot\frac13-\frac12\left(2-\left(-\frac14\right)\right)=1+1-\frac12\cdot\frac94=2-\frac98=\frac{16-9}{8}=\frac78.
\]
Torej je iskana vrednost $\dfrac78$.

\textbf{b)} $\log x=2\log3+\dfrac12\log9-\dfrac{1}{\log_{27}10}$. Ker je $\log_{27}10=\dfrac{\log10}{\log27}=\dfrac{1}{\log27}$ (saj je $\log10=1$), sledi
\[
\frac{1}{\log_{27}10}=\log27.
\]
Nadalje je $2\log3=\log9$, $\dfrac12\log9=\dfrac12\log(3^2)=\log3$ in $\log27=\log(3^3)=3\log3$. Torej
\[
\log x=\log9+\log3-\log27=\log\!\left(\frac{9\cdot3}{27}\right)=\log1=0\ \Longrightarrow\ x=1.
\]
\end{resitev}

\begin{resitev}[-- logaritemska ena\v{c}ba s produktom re\v{s}itev]
Definicijsko obmo\v{c}je: $x>0$. Levo stran preoblikujemo s pravilom $\log_b(a^c)=c\log_b a$:
\[
\log_2\!\left(x^{\log_2 x}\right)=(\log_2 x)\cdot(\log_2 x)=(\log_2 x)^2.
\]
Desno stran zdru\v{z}imo:
\[
\log_2(x^2)+\log_2 4=2\log_2 x+2.
\]
Z uvedbo $t=\log_2 x$ dobimo kvadratno ena\v{c}bo
\[
t^2=2t+2\ \Longrightarrow\ t^2-2t-2=0.
\]
Re\v{s}itvi sta
\[
t_{1,2}=\frac{2\pm\sqrt{4+8}}{2}=1\pm\sqrt3.
\]
Ker je $t=\log_2 x$, sta re\v{s}itvi za $x$ enaki
\[
x_1=2^{\,1+\sqrt3},\qquad x_2=2^{\,1-\sqrt3},
\]
obe pa ustrezata definicijskemu obmo\v{c}ju ($x_1,x_2>0$). Produkt re\v{s}itev je
\[
x_1\cdot x_2=2^{\,(1+\sqrt3)+(1-\sqrt3)}=2^{2}=\boldsymbol{4}.
\]
\end{resitev}

\end{document}
